\documentclass[11pt,a4paper]{article}
\usepackage[T1]{fontenc}
\usepackage[utf8]{inputenc}
\usepackage{amsmath,amsthm,mathtools}
\usepackage{newtxtext,newtxmath}
\usepackage[margin=27mm,headheight=15pt]{geometry}
\usepackage{microtype}
\usepackage{booktabs,array,longtable}
\usepackage[table]{xcolor}
\usepackage{enumitem}
\usepackage{titlesec,fancyhdr}
\usepackage[most]{tcolorbox}
\usepackage{xurl}
\usepackage{listings}
\usepackage{hyperref}
\definecolor{Olive}{HTML}{536348}
\definecolor{PaleOlive}{HTML}{F2F4ED}
\definecolor{Ink}{HTML}{262B25}
\definecolor{Muted}{HTML}{60665E}
\hypersetup{colorlinks=true,linkcolor=Olive,citecolor=Olive,urlcolor=Olive,
  pdftitle={Delayed digit stabilization: a counterexample and sharp bounds},
  pdfauthor={Research note prepared for Vladimir Reshetnikov},
  pdfsubject={An open stabilization conjecture, exact counterexamples, valuations, and fast-growing powers}}
\titleformat{\section}{\Large\bfseries\color{Olive}}{\thesection}{0.65em}{}
\titleformat{\subsection}{\large\bfseries\color{Olive}}{\thesubsection}{0.65em}{}
\titleformat{\subsubsection}{\normalsize\bfseries}{\thesubsubsection}{0.65em}{}
\pagestyle{fancy}\fancyhf{}
\fancyhead[L]{\small\color{Muted}Delayed digit stabilization}
\fancyhead[R]{\small\color{Muted}Research note \textbullet\ September 2026}
\fancyfoot[C]{\thepage}
\renewcommand{\headrulewidth}{0.3pt}
\setlength{\parindent}{1.2em}
\setlength{\parskip}{3pt}
\setlength{\emergencystretch}{2.5em}
\setlist{nosep,leftmargin=1.8em}
\newtheorem{theorem}{Theorem}[section]
\newtheorem{lemma}[theorem]{Lemma}
\newtheorem{proposition}[theorem]{Proposition}
\newtheorem{corollary}[theorem]{Corollary}
\theoremstyle{definition}
\newtheorem{definition}[theorem]{Definition}
\newtheorem{example}[theorem]{Example}
\newtheorem{remark}[theorem]{Remark}
\newcommand{\vp}[2]{\nu_{#1}\!\left(#2\right)}
\newcommand{\Z}{\mathbb Z}
\newcommand{\N}{\mathbb N}
\newcommand{\ord}{\operatorname{ord}}
\newcommand{\rad}{\operatorname{rad}}
\newcommand{\gcdop}{\operatorname{gcd}}
\newcommand{\lcm}{\operatorname{lcm}}
\newcommand{\astar}{a_{\star}}
\newtcolorbox{resultbox}[1][]{colback=PaleOlive,colframe=Olive,
  boxrule=0.5pt,arc=1.2mm,left=3mm,right=3mm,top=2mm,bottom=2mm,breakable,#1}
\lstset{basicstyle=\ttfamily\small,columns=fullflexible,breaklines=true,
  frame=single,rulecolor=\color{Olive},backgroundcolor=\color{PaleOlive},
  showstringspaces=false,aboveskip=9pt,belowskip=9pt}

\begin{document}
\begin{center}
{\small\sffamily\color{Olive}MODULAR ARITHMETIC OF RAPIDLY GROWING POWERS}\par\vspace{12pt}
{\LARGE\bfseries Delayed digit stabilization}\par\vspace{6pt}
{\Large A counterexample, exact minima, and sharp bounds}\par\vspace{14pt}
{\normalsize Research note prepared for Vladimir Reshetnikov}\par\vspace{4pt}
{\normalsize 19 September 2026}
\end{center}
\vspace{6pt}
\begin{abstract}
For $a\ge2$ not divisible by $10$, put
$S_a(b)=\nu_{10}(a^{10^{b+1}}-a^{10^b})$ and
$D_a(b)=S_a(b)-S_a(b-1)$. A publicly posed conjecture asks whether
$D_a(b)=1$ always holds for $b\ge3$. We disprove it and prove that the
smallest counterexample is
$\astar=3\cdot2^{99}+1=1901475900342344102245054808065$:
$S_{\astar}(2)=100$ and $S_{\astar}(3)=102$.
The mechanism is competition between an exponentially growing valuation at
one prime and an affine valuation at the other. We give exact formulas for
all bases, determine the first permanently unit increment, construct
arbitrarily delayed examples, and determine the least counterexample at
every prescribed stage. An exact density formula quantifies the scarcity of
exceptions. We then prove an arbitrary-radix local formula, characterize
convergence and finite limit cycles, and explain how fast-growing changes
of index transform the stabilization rate. A separate calculation for genuine
integer tetration makes the distinction between the two iterations explicit.
All mathematical claims are proved in ordinary mathematics; accompanying
standard-library Python code supplies independent modular certificates and
regression tests.
\end{abstract}
\begin{resultbox}
\textbf{Outcome.} The uniform threshold is false, not merely unproved.
Eventual unit speed for each fixed admissible base is true. The exact
threshold is unbounded as the base varies, even though it grows only on a
$\log\log a$ scale in the worst case.
\end{resultbox}
\vspace{4pt}
\noindent\textbf{Keywords.} Trailing digits; lifting the exponent; tetration;
$p$-adic valuation; finite iteration; Chinese remainder theorem.\par
\noindent\textbf{Scope.} The conjecture concerns $a^{10^b}$, not $a\uparrow\uparrow b$.
The relation to tetration and higher-growth indexing is developed in
Sections~\ref{sec:time} and~\ref{sec:tetration}.
\newpage
\begingroup
\setlength{\parskip}{0pt}
\small
\tableofcontents
\endgroup
\newpage

\section{The selected problem and its status}\label{sec:status}

\subsection{A concrete, independently checkable target}
The target is a question posed by Marco Rip\`a on Mathematics Stack Exchange
on 24 October 2025~\cite{question}. With the notation used in the abstract,
it asks whether every integer $a>1$ with $10\nmid a$ satisfies
\begin{equation}\label{eq:conjecture}
D_a(b)=1\qquad\text{for every }b\ge3.
\end{equation}
The source connects this question to earlier work on the congruence speed
of integer tetration. Its retrieved public page displayed no answers during
the literature check for this note. A targeted search did not locate a
resolution of this exact question. That is a record of the checked sources,
not a claim that every possible prior publication has been excluded.

The advantage of this target is that the conjectural assertion has precise
quantifiers and admits a finite certificate of failure. There is no need to
choose an analytic extension of tetration or a convention for ordinal
fundamental sequences. More importantly, the question can be reduced to a
competition between two exactly computable integer quantities.

\subsection{What is proved, and what is inherited}
The proof of the counterexample is in Section~\ref{sec:counterexample}; its
minimality is also proved there. The stronger least-counterexample theorem
for every stage appears in Section~\ref{sec:minima}. The density calculation
and the arbitrary-radix results are additional deductions developed here.
Their proofs, rather than an assertion of novelty, are the basis for the
claims made in this note.

The lifting identities used below are classical. The question already reports
a base-dependent sufficient onset bound~\cite{question}; pointwise eventual
unit speed is not claimed here as new. Earlier tetration papers
establish valuation formulas and eventual constant congruence speeds for
integer towers~\cite{speed,stable}. Those results are relevant background,
but they must not be substituted for the formulas of this note: the exponent
schedule is different. Work on the stable digits of Graham's number provides
another explicit application of valuation methods to very large
integers~\cite{graham}. We do not claim a new solution of that problem.

All calculations here use positive integer bases. Bases divisible by $10$
are included in the general formulas but excluded from the conjecture.
The base $1$ is excluded because every difference is zero and its valuation
would be infinite.

\subsection{The mathematical mechanism in one paragraph}
An even base can be exceptionally close to a fourth root of unity in the
$5$-adic metric while being divisible by exactly one power of $2$.
Conversely, an odd multiple of $5$ can be exceptionally close to $1$ or $-1$
in the $2$-adic metric while containing only one factor of $5$.
The number of matching decimal digits is the \emph{smaller} of the two local
valuations. A large constant in the affine local valuation can therefore
postpone the point at which it becomes the smaller one. This is the missing
uniformity in~\eqref{eq:conjecture}.

\section{Valuations and the lifting identities}\label{sec:lte}

For a prime $p$ and a nonzero integer $z$, let $\nu_p(z)$ be the largest
integer $k\ge0$ such that $p^k\mid z$. For any radix $B\ge2$, define
\begin{equation}\label{eq:radix-valuation}
\nu_B(z)=\max\{k\ge0:B^k\mid z\}.
\end{equation}
If $B=\prod_{p\mid B}p^{e_p}$, then
\begin{equation}\label{eq:minval}
\nu_B(z)=\min_{p\mid B}\left\lfloor\frac{\nu_p(z)}{e_p}\right\rfloor.
\end{equation}
This follows by comparing the prime exponents in $B^k$ and $z$.
In particular, $\nu_{10}(z)=\min\{\nu_2(z),\nu_5(z)\}$.
For a composite radix this is a divisibility count; it is not generally an
additive prime valuation.

Two nonnegative integers agree in their last $k$ base-$B$ positions exactly
when their difference is divisible by $B^k$. Leading zero padding is
understood when necessary.

\begin{lemma}[Odd-prime lifting]\label{lem:odd-lte}
Let $p$ be odd and let $p\mid u-1$. For every positive integer $m$,
\[
\nu_p(u^m-1)=\nu_p(u-1)+\nu_p(m).
\]
More generally, if $p\nmid a$ and $d=\ord_p(a)$, then
\[
\nu_p(a^m-1)=
\begin{cases}
0,&d\nmid m,\\
\nu_p(a^d-1)+\nu_p(m),&d\mid m.
\end{cases}
\]
\end{lemma}
\begin{proof}
Write $m=p^j w$, where $p\nmid w$. If $u\equiv1\pmod p$, then
$(u^w-1)/(u-1)$ is congruent to $w$ modulo $p$, so raising to the power $w$
does not change the valuation. If $u=1+t$ with $\nu_p(t)=r\ge1$, the first
term of
\[
(1+t)^p-1=pt+\binom p2t^2+\cdots+t^p
\]
has valuation $r+1$. Every other term has valuation at least $r+2$.
Thus each subsequent raising to the power $p$ increases the valuation by
exactly one. This proves the first identity.

For the second statement, $a^m\equiv1\pmod p$ holds exactly when $d\mid m$.
In that case apply the first identity to $u=a^d$ and exponent $m/d$.
Since $d\mid p-1$, we have $p\nmid d$ and $\nu_p(m/d)=\nu_p(m)$.
\end{proof}

\begin{lemma}[Binary lifting]\label{lem:two-lte}
For odd $a$ and positive even $m$,
\begin{equation}\label{eq:binary-lte}
\nu_2(a^m-1)=\nu_2(a-1)+\nu_2(a+1)+\nu_2(m)-1.
\end{equation}
\end{lemma}
\begin{proof}
Write $m=2^j w$ with $j\ge1$ and $w$ odd. Odd powers leave the valuation of
$u-1$ unchanged, since $(u^w-1)/(u-1)$ is an odd sum of odd terms.
The first squaring contributes
$\nu_2(a^2-1)=\nu_2(a-1)+\nu_2(a+1)$, which is at least $3$.
After this squaring the current value is $1$ modulo $8$, so each further
factor $u+1$ contributes exactly one factor of $2$. There are $j-1$ such
further squarings, giving~\eqref{eq:binary-lte}.
\end{proof}

\begin{lemma}[A prime dividing the base]\label{lem:nonunit}
If $p\mid a$ and $0<u<v$, then
\[
\nu_p(a^v-a^u)=u\nu_p(a).
\]
\end{lemma}
\begin{proof}
Factor the difference as $a^u(a^{v-u}-1)$. The second factor is $-1$ modulo
$p$ and contributes no further factors of $p$.
\end{proof}

\section{The complete decimal profile}\label{sec:profile}

For an odd integer $a>1$, put
\begin{equation}\label{eq:c2}
C_2(a)=\nu_2(a^2-1)-1
      =\max\{\nu_2(a-1),\nu_2(a+1)\}.
\end{equation}
The equality holds because exactly one of $a-1,a+1$ is divisible by $4$,
and the other has valuation $1$. For $5\nmid a$, put
\begin{equation}\label{eq:c5}
C_5(a)=\nu_5(a^4-1).
\end{equation}
The factorization of $a^4-1$ gives the useful equivalent formula
\begin{equation}\label{eq:c5-split}
C_5(a)=
\begin{cases}
\nu_5(a-1),&a\equiv1\pmod5,\\
\nu_5(a+1),&a\equiv4\pmod5,\\
\nu_5(a^2+1),&a\equiv2,3\pmod5.
\end{cases}
\end{equation}
Only one of the displayed factors is divisible by $5$.

\begin{theorem}[Exact profile for every decimal base]\label{thm:profile}
Let $a\ge2$ and $b\ge2$. Then
\begin{equation}\label{eq:profile}
S_a(b)=
\begin{cases}
b+\min\{C_2(a),C_5(a)\},&\gcd(a,10)=1,\\[2pt]
\min\{\nu_2(a)10^b,\ b+C_5(a)\},&2\mid a,\ 5\nmid a,\\[2pt]
\min\{\nu_5(a)10^b,\ b+C_2(a)\},&5\mid a,\ 2\nmid a,\\[2pt]
10^b\min\{\nu_2(a),\nu_5(a)\},&10\mid a.
\end{cases}
\end{equation}
\end{theorem}
\begin{proof}
The fundamental factorization is
\begin{equation}\label{eq:factorization}
a^{10^{b+1}}-a^{10^b}=a^{10^b}\bigl(a^{9\cdot10^b}-1\bigr).
\end{equation}
If $2\mid a$, Lemma~\ref{lem:nonunit} gives valuation $\nu_2(a)10^b$.
If $a$ is odd, Lemma~\ref{lem:two-lte} applies with
$\nu_2(9\cdot10^b)=b$ and gives $b+C_2(a)$.

If $5\mid a$, Lemma~\ref{lem:nonunit} gives valuation $\nu_5(a)10^b$.
If $5\nmid a$, then $\ord_5(a)\mid4$, while $4\mid10^b$ because $b\ge2$.
Apply odd-prime lifting to $a^4$ with exponent $9\cdot10^b/4$.
The result is $b+C_5(a)$.
Taking the smaller of the binary and quinary valuations proves all four
cases.
\end{proof}

The hypothesis $b\ge2$ is substantive: when $b=1$, the exponent
$9\cdot10=90$ need not be divisible by $4$. The conjectured increments
start at $b=3$, so both $S_a(b)$ and $S_a(b-1)$ lie in the range where the
theorem applies.

\begin{corollary}[The two different assertions]\label{cor:fixed}
For every fixed $a\ge2$ with $10\nmid a$, there is an integer $C(a)$ such
that $S_a(b)=b+C(a)$ for all sufficiently large $b$.
If $\gcd(a,10)=1$, this already holds for every $b\ge2$.
\end{corollary}
\begin{proof}
The coprime case is explicit in~\eqref{eq:profile}. In either mixed case,
$t10^b>b+C$ eventually, for the fixed positive integers $t$ and $C$.
\end{proof}
This proves pointwise eventual stabilization. It gives no base-independent
bound on the initial exceptional interval.

\section{A counterexample and the proof of its minimality}\label{sec:counterexample}

\subsection{The counterexample}
\begin{theorem}\label{thm:counterexample}
The integer
\begin{equation}\label{eq:astar}
\astar=3\cdot2^{99}+1=1901475900342344102245054808065
\end{equation}
is a counterexample to~\eqref{eq:conjecture}. Specifically,
\[
S_{\astar}(2)=100,\qquad S_{\astar}(3)=102,\qquad D_{\astar}(3)=2.
\]
\end{theorem}
\begin{proof}
We have $2^{99}\equiv13\pmod{25}$, so $\astar\equiv15\pmod{25}$.
Consequently $\nu_5(\astar)=1$. Moreover,
\[
\nu_2(\astar-1)=99,\qquad \nu_2(\astar+1)=1,
\]
and hence $C_2(\astar)=99$. The odd-multiple-of-$5$ case of
Theorem~\ref{thm:profile} now gives
\begin{equation}\label{eq:astar-profile}
S_{\astar}(b)=\min\{10^b,b+99\}\qquad(b\ge2).
\end{equation}
At $b=2$ the minimum is $100$, whereas at $b=3$ it is $102$.
\end{proof}

\begin{center}
\begin{tabular}{rrrrr}
\toprule
$b$ & $\nu_2(\Delta_b)$ & $\nu_5(\Delta_b)$ & $S_{\astar}(b)$ & $D_{\astar}(b)$\\
\midrule
2 & 101 & 100 & 100 & ---\\
3 & 102 & 1000 & 102 & 2\\
4 & 103 & 10000 & 103 & 1\\
5 & 104 & 100000 & 104 & 1\\
6 & 105 & 1000000 & 105 & 1\\
\bottomrule
\end{tabular}\par\smallskip
{\small Here $\Delta_b=\astar^{10^{b+1}}-\astar^{10^b}$.}
\end{center}
The exceptional increment is caused by a change in which prime controls
the minimum. No approximation of the enormous powers is involved.

\subsection{An exact test for unit increments}
\begin{lemma}\label{lem:increment}
Let $t,C$ be positive integers, and let
$S(b)=\min\{t10^b,b+C\}$ for $b\ge2$. For $b\ge3$,
\begin{equation}\label{eq:criterion}
S(b)-S(b-1)=1
\quad\Longleftrightarrow\quad
C\le t10^{b-1}-(b-1).
\end{equation}
If the inequality fails, the increment is at least $2$.
\end{lemma}
\begin{proof}
Write $R=t10^{b-1}$ and $L=b-1+C$. If $L\le R$, then the previous
minimum is $L$ and the next is $L+1$, because $L+1\le R+1\le10R$.
If $L>R$, integrality gives $L\ge R+1$. The previous minimum is $R$ and
the next is $\min\{10R,L+1\}\ge R+2$.
\end{proof}

Since $t10^{b-1}-(b-1)$ strictly increases with $b$, once a unit increment
occurs, every later increment is also one. In particular, a base satisfies
the conjecture at every stage $b\ge3$ exactly when it satisfies it at $b=3$.

\subsection{Why no smaller integer works}
\begin{theorem}[Global minimality]\label{thm:minimal}
The integer $\astar$ of~\eqref{eq:astar} is the smallest $a\ge2$, $10\nmid a$,
for which $D_a(b)\ne1$ at any stage $b\ge3$.
\end{theorem}
\begin{proof}
It suffices to rule out $D_a(3)\ne1$ for $a<\astar$.
Coprime bases cannot fail, by Theorem~\ref{thm:profile}.

Suppose first that $a$ is even and $5\nmid a$. By Lemma~\ref{lem:increment},
failure requires
\[
C_5(a)\ge100\nu_2(a)-1\ge99.
\]
By~\eqref{eq:c5-split}, one of $a-1,a+1,a^2+1$ is divisible by $5^{99}$.
Thus $a^2+1\ge5^{99}$. But
\begin{equation}\label{eq:small-bound}
\astar^2+1<5^{99},
\end{equation}
so this is impossible when $a<\astar$. For an elementary check of the
bound, $\astar<4\cdot2^{99}$, and
$16\cdot4^{99}+1<17\cdot4^{99}<5^{99}$.
The latter inequality follows already from $(5/4)^{16}>25>17$ and $99>16$.

Now suppose that $a$ is odd and $5\mid a$. Failure requires
\[
C_2(a)\ge100\nu_5(a)-1\ge99.
\]
Hence $a\equiv1$ or $-1\pmod{2^{99}}$. Among the integers between $2$ and
$3\cdot2^{99}+1$, the only possible smaller multiple of $5$ in these two
residue classes is $2^{100}-1$. This can be checked by considering
$j2^{99}\pm1$ for $j=1,2,3$ and using $2^{99}\equiv3\pmod5$.
For that remaining candidate,
\[
C_2(2^{100}-1)=100,
\qquad
\nu_5(2^{100}-1)=\nu_5(16^{25}-1)=1+2=3.
\]
Its failure condition would be $100\ge299$, which is false.
Thus every smaller admissible base satisfies the conjectured bound.
\end{proof}

The size of the first counterexample explains why a search restricted to
small bases offers misleading evidence for the uniform statement. This is
a conclusion from a minimality proof, not an extrapolation of an experiment.

\section{The exact onset and arbitrarily long delays}\label{sec:onset}

Define the onset, restricted to the range relevant to the conjecture, by
\begin{equation}\label{eq:onset-def}
N(a)=\min\{B\ge3:D_a(b)=1\text{ for all }b\ge B\}.
\end{equation}
It exists for every admissible $a$ by Corollary~\ref{cor:fixed}.

\begin{theorem}[Exact onset]\label{thm:onset}
If $\gcd(a,10)=1$, then $N(a)=3$. In either mixed case, use
\[
(t,C)=
\begin{cases}
(\nu_2(a),C_5(a)),&2\mid a,\ 5\nmid a,\\
(\nu_5(a),C_2(a)),&5\mid a,\ 2\nmid a.
\end{cases}
\]
Then
\begin{equation}\label{eq:onset}
N(a)=1+\min\{r\ge2:t10^r-r\ge C\}.
\end{equation}
\end{theorem}
\begin{proof}
Substitute $r=b-1$ in Lemma~\ref{lem:increment}; the inequality is
monotone in $r$.
\end{proof}

This is a sharp characterization, not merely a sufficient bound. If
\[
r_0=\max\left\{2,\left\lceil\log_{10}(C/t)\right\rceil\right\},
\]
the minimum in~\eqref{eq:onset} is either $r_0$ or $r_0+1$.
Indeed, smaller admissible $r$ cannot make $t10^r\ge C$, and the additional
factor of ten at the next stage more than compensates for subtracting $r$.
The exact formula can be evaluated by integer comparisons without computing
any logarithms.

\begin{theorem}[A family realizing every later onset]\label{thm:delay}
For $m\ge2$, set
\begin{equation}\label{eq:delay-family}
K_m=10^m-m+1,\qquad a_m=5^{K_m}+1.
\end{equation}
Then
\[
D_{a_m}(m+1)=2,\qquad N(a_m)=m+2.
\]
In particular, no absolute onset bound is valid for all admissible bases.
\end{theorem}
\begin{proof}
For every positive integer $K$, the number $5^K+1$ is $2$ modulo $4$ and
$1$ modulo $5$. Thus $\nu_2(5^K+1)=1$ and $C_5(5^K+1)=K$, so
\[
S_{5^K+1}(b)=\min\{10^b,b+K\}.
\]
At $K=K_m$, the linear term exceeds $10^m$ by one when $b=m$.
Hence $S_{a_m}(m)=10^m$, while $S_{a_m}(m+1)=10^m+2$.
The criterion fails at $r=m$ and holds at $r=m+1$; because it is monotone,
the asserted onset is exact.
\end{proof}

The family $5^K+1$ gives $N(5^K+1)=\log_{10}K+O(1)$ as $K\to\infty$.
More generally,
\begin{equation}\label{eq:loglog}
N(a)=O(\log\log a)
\end{equation}
as $a$ ranges over admissible bases. To see this, use
$C_2(a)\le\log_2(a+1)$ and $C_5(a)\le\log_5(a^2+1)$ in the onset
formula, together with $t\ge1$. The family above also shows that the
worst-case order in~\eqref{eq:loglog} cannot be improved to $o(\log\log a)$.
Unboundedness and slow growth are entirely compatible.

\section{The least counterexample at every stage}\label{sec:minima}

The preceding family is convenient but not usually minimal. We can solve
the minimization problem exactly.

\begin{definition}
For $b\ge3$, let
\[
A_b=\min\{a\ge2:10\nmid a,\ D_a(b)\ne1\},
\qquad K=10^{b-1}-b+2.
\]
\end{definition}
The set is nonempty by Theorem~\ref{thm:delay}.

\begin{theorem}[Exact stage minima]\label{thm:stage-minima}
For $b\ge3$, the least exceptional base is the following integer:
\begin{center}
\begin{tabular}{lll}
\toprule
Condition on $K$ & $A_b$ & $D_{A_b}(b)$\\
\midrule
$K\equiv0\pmod4$, $K\not\equiv0\pmod{20}$ & $2^K-1$ & 2\\
$K\equiv0\pmod{20}$ & $4\cdot2^K+1$ & 4\\
$K\equiv1\pmod4$, $K\not\equiv9\pmod{20}$ & $2\cdot2^K+1$ & 3\\
$K\equiv9\pmod{20}$ & $3\cdot2^K-1$ & 2\\
$K\equiv2\pmod4$, $K\not\equiv10\pmod{20}$ & $2^K+1$ & 2\\
$K\equiv10\pmod{20}$ & $4\cdot2^K-1$ & 4\\
$K\equiv3\pmod4$, $K\not\equiv19\pmod{20}$ & $2\cdot2^K-1$ & 3\\
$K\equiv19\pmod{20}$ & $3\cdot2^K+1$ & 2\\
\bottomrule
\end{tabular}
\end{center}
In every case $N(A_b)=b+1$.
\end{theorem}
\begin{proof}
Put $r=b-1$ and $Q=10^r$, so $K=Q-r+1\ge99$.
For a mixed base, the failure condition is
\begin{equation}\label{eq:stage-failure}
C\ge tQ-r+1.
\end{equation}
We will find the minimum below $4\cdot2^K+1$.

An even exceptional base must have $C_5(a)\ge K$, so
$a^2+1\ge5^K$. But for $K\ge99$,
\[
(4\cdot2^K+1)^2+1<25\cdot4^K<5^K.
\]
The first inequality holds already for $K\ge1$; the second follows from
$(5/4)^{16}>25$ and $K\ge99$. No even base in the proposed range can
therefore be exceptional.

For an odd multiple of $5$ in this range, $C_2(a)\le K+2$.
If $\nu_5(a)\ge2$, condition~\eqref{eq:stage-failure} requires
\[
C_2(a)\ge2Q-r+1=2K+r-1>K+2,
\]
which is impossible. Thus we need precisely
\begin{equation}\label{eq:min-conditions}
\nu_5(a)=1,\qquad a\equiv\pm1\pmod{2^K}.
\end{equation}
Every candidate in the range, other than the excluded integer $1$, is
$j2^K\pm1$ for $1\le j\le4$.

The residue of $2^K$ modulo $5$ gives the first possible multiple of $5$.
For $K$ congruent to $0,1,2,3$ modulo $4$, respectively, these first
candidates are
\[
2^K-1,\qquad 2\cdot2^K+1,\qquad
2^K+1,\qquad 2\cdot2^K-1.
\]
A candidate must be skipped exactly when it is also divisible by $25$.
The order of $2$ modulo $25$ is $20$, with $2^{10}\equiv-1\pmod{25}$.
The four exceptional residue classes are consequently
$K\equiv0,9,10,19\pmod{20}$, respectively. The next candidates are
$4\cdot2^K+1$, $3\cdot2^K-1$, $4\cdot2^K-1$, and $3\cdot2^K+1$.
Direct reduction modulo $25$ shows that each has exactly one factor of $5$.
This proves the table and its minimality.

If the selected candidate is $A_b=j2^K\pm1$, then
$C_2(A_b)=K+\nu_2(j)$. At stage $b-1$, its profile is $Q$; at stage $b$,
it is $b+K+\nu_2(j)$. Their difference is
$2+\nu_2(j)$, giving the last column.
At the next stage the inequality of Lemma~\ref{lem:increment} holds, so
$N(A_b)=b+1$.
\end{proof}

Because $10^{b-1}\equiv0\pmod{20}$, the coefficient and sign in the table
depend only on $b$ modulo $20$: indeed $K\equiv2-b\pmod{20}$.
The first few values are
\begin{center}
\begin{tabular}{rll}
\toprule
$b$ & $K$ & $A_b$\\
\midrule
3 & 99 & $3\cdot2^{99}+1$\\
4 & 998 & $2^{998}+1$\\
5 & 9997 & $2^{9998}+1$\\
6 & 99996 & $2^{99996}-1$\\
7 & 999995 & $2^{999996}-1$\\
\bottomrule
\end{tabular}
\end{center}

\begin{corollary}[An exact extremal description]
For $X\ge2$, let $M(X)$ be the maximum of $N(a)$ over admissible integers
$a\le X$. If $X<A_3$, then $M(X)=3$. Otherwise,
\[
M(X)=1+\max\{b\ge3:A_b\le X\}.
\]
In particular,
\[
\log_2 A_b=10^{b-1}-b+O(1),\qquad
M(X)=\log_{10}\log_2 X+O(1)\quad(X\to\infty).
\]
\end{corollary}
\begin{proof}
The statement $N(a)>b$ is equivalent to $D_a(b)\ne1$.
The definition of $A_b$ therefore characterizes exactly when the set
$a\le X$ contains a base with onset exceeding $b$. The asymptotic estimates
follow from $A_b=j2^K\pm1$ with $1\le j\le4$.
\end{proof}

\section{How sparse are the counterexamples?}\label{sec:density}

For a fixed $b\ge3$, write
\[
E_b=\{a\ge2:10\nmid a,\ D_a(b)\ne1\}.
\]
We compute its natural density among \emph{all} positive integers:
\[
\delta_b=\lim_{X\to\infty}\frac{\#(E_b\cap[1,X])}{X}.
\]
The conditional density among admissible bases is $10\delta_b/9$.

\begin{theorem}[Exact exceptional density]\label{thm:density}
Let $r=b-1$ and $Q=10^r$. The density exists and equals
\begin{equation}\label{eq:density}
\boxed{\displaystyle
\delta_b=
\frac{2\cdot5^{r-1}}{2\cdot5^Q-1}
+\frac{2^{r+2}}{5(5\cdot2^Q-1)}.}
\end{equation}
The first term counts even bases; the second counts odd multiples of $5$.
\end{theorem}
\begin{proof}
Fix an even base with $\nu_2(a)=t\ge1$. Failure means
\[
\nu_5(a^4-1)\ge K_t,\qquad K_t=tQ-r+1.
\]
There are exactly four roots of $x^4=1$ modulo $5^{K_t}$.
Indeed, modulo $5$ all four nonzero classes are roots, and the derivative
$4x^3$ is nonzero at each root. To see the unique lifting directly, replace
a root $x$ modulo $5^k$ by $x+j5^k$; the coefficient of $j$ in the next
congruence is invertible modulo $5$, determining exactly one $j$.
The density of the quinary condition is thus $4/5^{K_t}$.
The exact binary valuation condition has density $1/2^{t+1}$.
By the Chinese remainder theorem their intersection has density
\[
\frac{4}{2^{t+1}5^{tQ-r+1}}.
\]
Summing over $t\ge1$ gives
\[
2\cdot5^{r-1}\sum_{t\ge1}\left(\frac1{2\cdot5^Q}\right)^t
=\frac{2\cdot5^{r-1}}{2\cdot5^Q-1}.
\]

For odd multiples of $5$ with $\nu_5(a)=s\ge1$, the density of the exact
quinary condition is $4/5^{s+1}$. Failure requires
\[
C_2(a)\ge K_s=sQ-r+1.
\]
Since $K_s\ge99$, this is exactly $a\equiv1$ or $-1\pmod{2^{K_s}}$,
a condition of density $2/2^{K_s}$. Again applying the Chinese remainder
theorem gives density
\[
\frac{8}{5^{s+1}2^{sQ-r+1}}.
\]
Its sum over $s\ge1$ is
\[
\frac{2^{r+2}}5\sum_{s\ge1}\left(\frac1{5\cdot2^Q}\right)^s
=\frac{2^{r+2}}{5(5\cdot2^Q-1)}.
\]

It remains to justify summing infinitely many densities. After truncating
the even-base sum at $t<T$, its omitted set is contained in the multiples
of $2^T$, whose upper density is $2^{-T}$. The corresponding omitted
odd-base set is contained in the multiples of $5^T$, of upper density
$5^{-T}$. These bounds tend to zero. Finite sums have densities by the
Chinese remainder theorem, so the infinite union has the stated density.
\end{proof}

For the original conjecture,
\begin{equation}\label{eq:density3}
\delta_3=
\frac{10}{2\cdot5^{100}-1}
+\frac{16}{5(5\cdot2^{100}-1)}
\approx5.0487097934144756\times10^{-31}.
\end{equation}
This is a natural density, not an independence assumption about random
samples. In particular, it should not be interpreted as a theorem about
the distribution of exceptions in every finite interval.

The second term dominates as $b\to\infty$, so
\begin{equation}\label{eq:density-asymptotic}
\delta_b\sim\frac{2^{b+1-10^{b-1}}}{25}.
\end{equation}
The extraordinary rarity of exceptions is therefore quantified by a
closed formula, while their existence at every stage is guaranteed by the
explicit constructions.

\section{Permanent digits and the four decimal limits}\label{sec:limits}

So far $S_a(b)$ was defined using neighboring terms. The same number actually
measures agreement with \emph{every} later term.

\begin{theorem}[Exact nonconsecutive agreement]\label{thm:nonconsecutive}
For every $a\ge2$ and integers $v>u\ge2$,
\begin{equation}\label{eq:nonconsecutive}
\nu_{10}\!\left(a^{10^v}-a^{10^u}\right)=S_a(u).
\end{equation}
\end{theorem}
\begin{proof}
Factor the difference as
\[
a^{10^u}\left(a^{10^u(10^{v-u}-1)}-1\right).
\]
The number $10^{v-u}-1$ is coprime to $10$. Consequently the exponent in
the second factor has binary and quinary valuations both equal to $u$,
and is divisible by $4$. Every step of the proof of
Theorem~\ref{thm:profile} therefore gives the same local valuations as for
the neighboring-term difference.
\end{proof}

The $10$-adic integers may be defined as compatible systems of residues
modulo $10^k$. The Chinese remainder theorem gives
\[
\Z_{10}\cong\Z_2\times\Z_5.
\]
This is a product of two rings, not a field. Since $S_a(u)\to\infty$,
Theorem~\ref{thm:nonconsecutive} shows that $a^{10^u}$ has a limit
$L_a\in\Z_{10}$.

\begin{theorem}[Limit and exact distance]\label{thm:idempotents}
The limit depends only on which of $2$ and $5$ divide $a$:
\begin{equation}\label{eq:limit-components}
L_a\longleftrightarrow
\bigl(\mathbf1_{2\nmid a},\mathbf1_{5\nmid a}\bigr)
\quad\text{in }\Z_2\times\Z_5.
\end{equation}
These are the four idempotents of $\Z_{10}$. Moreover, for $u\ge2$,
\begin{equation}\label{eq:distance}
\nu_{10}(L_a-a^{10^u})=S_a(u).
\end{equation}
\end{theorem}
\begin{proof}
If $p\mid a$, then $\nu_p(a^{10^u})=10^u\nu_p(a)\to\infty$, so the
$p$-adic component tends to zero. If $a$ is odd, binary lifting gives
$\nu_2(a^{10^u}-1)=u+C_2(a)\to\infty$.
If $5\nmid a$ and $u\ge2$, odd-prime lifting gives
$\nu_5(a^{10^u}-1)=u+C_5(a)\to\infty$.
This proves~\eqref{eq:limit-components}. In each $\Z_p$, the equation
$x(x-1)=0$ has only the solutions $0,1$, since $\Z_p$ has no zero divisors.
Thus the product ring has exactly the four displayed idempotents.

For the exact distance, $S_a(u+1)>S_a(u)$ in all cases of
Theorem~\ref{thm:profile}. Every later difference from $a^{10^{u+1}}$
is divisible by $10^{S_a(u)+1}$. Passing to the limit gives
\[
L_a-a^{10^u}\equiv a^{10^{u+1}}-a^{10^u}\pmod{10^{S_a(u)+1}}.
\]
The right side has valuation exactly $S_a(u)$, proving~\eqref{eq:distance}.
\end{proof}

The finite residue of $L_a$ can be computed without iteration. To find its
last $k$ digits, solve
\[
z\equiv\mathbf1_{2\nmid a}\pmod{2^k},\qquad
z\equiv\mathbf1_{5\nmid a}\pmod{5^k}
\]
by the Chinese remainder theorem. Thus the asymptotic limit is elementary;
the interesting issue is the exact rate of approach and its nonuniform
initial behavior.

\section{Fast-growing changes of index}\label{sec:time}

Let $g:\N\to\N$ be any strictly increasing function with $g(n)\ge2$, and
consider the sampled sequence
\[
Y_n=a^{10^{g(n)}}.
\]
No computability or growth-rate assumption is needed for the next statement.

\begin{theorem}[A time-change identity]\label{thm:time-change}
For every $n$,
\[
\nu_{10}(Y_{n+1}-Y_n)=S_a(g(n)).
\]
For fixed $a\ge2$ with $10\nmid a$, there is a fixed integer $C(a)$ such
that eventually
\begin{align}
\nu_{10}(Y_{n+1}-Y_n)&=g(n)+C(a),\label{eq:time-count}\\
\nu_{10}(Y_{n+1}-Y_n)-\nu_{10}(Y_n-Y_{n-1})
&=g(n)-g(n-1).\label{eq:time-speed}
\end{align}
\end{theorem}
\begin{proof}
Apply Theorem~\ref{thm:nonconsecutive} with $u=g(n)$ and $v=g(n+1)$.
A strictly increasing integer-valued $g$ is unbounded, so eventually
$g(n)$ lies in the affine part of the fixed-base profile.
Subtract the two consecutive affine formulas.
\end{proof}

For example, choose $g(n)=2\uparrow\uparrow n$ for $n\ge2$, where
$2\uparrow\uparrow n$ denotes the right-associated tower of $n$ copies of
$2$. Then the exact eventual count of new stable digits per step is
\[
2\uparrow\uparrow n-2\uparrow\uparrow(n-1).
\]
The same identity applies to any specified increasing integer-valued
schedule obtained from higher up-arrow expressions, a chain notation, or a
chosen fast-growing hierarchy. The theorem concerns sampling an already
defined sequence; it does not assert identities between those different
large-number notations.

This distinction matters conceptually. A stabilization rate is attached
not just to the limit and base but also to the chosen indexing scheme.
Passing from $n$ to a tetrational schedule can turn unit increments into
tetrationally large increments without changing the limiting $10$-adic
integer.

\section{An exact local formula in every radix}\label{sec:radix}

Now let $B\ge2$ be arbitrary, with
$B=\prod_{p\mid B}p^{e_p}$. For $a\ge2$, $n\ge1$, and a positive lag $h$,
put
\[
S_{a,B,h}(n)=\nu_B\!\left(a^{B^{n+h}}-a^{B^n}\right).
\]
When $h=1$ write $S_{a,B}(n)$.

For an odd prime $p\mid B$ with $p\nmid a$, define
\[
d_p=\ord_p(a),\qquad c_p=\nu_p(a^{d_p}-1).
\]
For $2\mid B$ and odd $a$, put $c_2=\nu_2(a^2-1)-1$.

\begin{theorem}[Arbitrary-radix profile]\label{thm:radix-profile}
For every $n,h\ge1$,
\begin{equation}\label{eq:radix-profile}
S_{a,B,h}(n)=\min_{p\mid B}\left\lfloor\frac{w_p(n,h)}{e_p}\right\rfloor,
\end{equation}
where
\begin{equation}\label{eq:wp}
w_p(n,h)=
\begin{cases}
\nu_p(a)B^n,&p\mid a,\\
n e_2+c_2,&p=2,\ 2\nmid a,\\
0,&p>2,\ p\nmid a,\ d_p\nmid B^n(B^h-1),\\
n e_p+c_p,&p>2,\ p\nmid a,\ d_p\mid B^n(B^h-1).
\end{cases}
\end{equation}
\end{theorem}
\begin{proof}
Factor the difference as
\[
a^{B^n}\left(a^{B^n(B^h-1)}-1\right).
\]
For $p\mid a$, use Lemma~\ref{lem:nonunit}. For a unit at an odd $p$,
use Lemma~\ref{lem:odd-lte}. Since $p\mid B$, the factor $B^h-1$ is
coprime to $p$, and the exponent has valuation $ne_p$.
At $p=2$, the exponent is even and the same observation applies to
Lemma~\ref{lem:two-lte}. Finally use~\eqref{eq:minval}.
\end{proof}

The multiplicative-order condition cannot be omitted. It is automatically
satisfied for decimal bases once $n\ge2$, but not for arbitrary radices.
The floors cannot be omitted either when $B$ contains prime powers.

\subsection{Convergence and finite limit cycles}
For an odd unit prime $p\mid B$, let $d_p^{\perp}$ be the largest divisor
of $d_p$ coprime to $B$: remove from $d_p$ all prime factors that divide $B$.

\begin{theorem}[Convergence criterion and cycle length]\label{thm:cycles}
The sequence $a^{B^n}$ converges in $\Z_B$ if and only if
\begin{equation}\label{eq:convergence}
d_p^{\perp}\mid B-1
\quad\text{for every odd }p\mid B\text{ with }p\nmid a.
\end{equation}
In general it has an asymptotic cycle whose exact length is
\begin{equation}\label{eq:cycle-length}
H=\lcm_{\substack{p\mid B,\ p>2\\p\nmid a}}
\ord_{d_p^{\perp}}(B),
\end{equation}
where an order modulo $1$, and an empty least common multiple, are taken
to be $1$. Precisely, each sequence with indices in one residue class
modulo $H$ converges, and $H$ is the least positive integer with that property.
If $H=1$, the limit $L$ satisfies $L^B=L$.
\end{theorem}
\begin{proof}
For sufficiently large $n$, every prime-power divisor of $d_p$ supported on
the primes dividing $B$ divides $B^n$. The remaining factor $d_p^{\perp}$
is coprime to $B^n$. Thus, eventually,
\[
d_p\mid B^n(B^h-1)
\quad\Longleftrightarrow\quad
 d_p^{\perp}\mid B^h-1.
\]
For $h=1$, if one condition fails, adjacent terms remain different modulo
that prime $p$, so the sequence cannot converge. If every condition holds,
Theorem~\ref{thm:radix-profile} makes every local valuation tend to infinity.
For each fixed precision all sufficiently late adjacent differences vanish;
then any finite sum of such differences vanishes at that precision. The
sequence is Cauchy and hence converges in the inverse-limit ring $\Z_B$.

The smallest positive $h$ satisfying all the displayed divisibilities is
exactly~\eqref{eq:cycle-length}. For that lag, the same valuation argument
makes every residue-class subsequence Cauchy. If a smaller positive lag
worked, its adjacent subsequence terms would eventually agree modulo every
$p\mid B$, forcing all the same divisibilities and contradicting the
minimality of $H$. Finally $X_{n+1}=X_n^B$ and continuity of polynomial
maps show that the limiting points are carried cyclically to one another;
in the convergent case this gives $L^B=L$.
\end{proof}

\begin{example}[An obstruction in radix $14$]
Take $a=2$, $B=14$. Modulo $7$, the order of $2$ is $3$, and
$3\nmid14-1$. Thus $S_{2,14}(n)=0$ for every $n\ge1$; adjacent terms never
share even one trailing radix-$14$ digit. Since $14\equiv-1\pmod3$, the
cycle length is $2$. For lag $2$ the local formulas give
\[
S_{2,14,2}(n)=\min\{14^n,n+1\}=n+1\qquad(n\ge1).
\]
The even and odd index subsequences therefore converge separately with an
exact affine agreement count.
\end{example}

\subsection{The eventual affine profile}
Assume~\eqref{eq:convergence}, and suppose at least one prime divisor of
$B$ does not divide $a$. Once the order conditions have become satisfied,
Theorem~\ref{thm:radix-profile} reduces to
\begin{equation}\label{eq:radix-envelope}
S_{a,B}(n)=\min\left\{n+C_B(a),\ 
\min_{\substack{p\mid B\\p\mid a}}
\left\lfloor\frac{\nu_p(a)B^n}{e_p}\right\rfloor\right\},
\end{equation}
where
\[
C_B(a)=\min_{\substack{p\mid B\\p\nmid a}}
\left\lfloor\frac{c_p}{e_p}\right\rfloor.
\]
The inner minimum over an empty set is omitted. Thus the eventual slope is
one whenever there is a unit component and convergence occurs.
If every prime dividing $B$ divides $a$, all components instead have
exponential profiles and the limit is zero.

For nonsquarefree $B$, the condition $B\nmid a$ does \emph{not} guarantee
a unit component. For instance $12\nmid6$, but both prime divisors of $12$
divide $6$. The correct condition is $\rad(B)\nmid a$.

\section{When can a radix have a uniform onset?}\label{sec:uniform-radix}

To compare radices fairly, restrict to bases for which $a^{B^n}$ converges
and at least one prime divisor of $B$ is a unit prime for $a$. These are the
bases for which the eventual unit-slope conclusion applies.

\begin{theorem}[Prime powers versus mixed radices]\label{thm:radix-dichotomy}
If $B$ is a prime power, every base in this class has
$S_{a,B}(n)=n+C_B(a)$ for all $n\ge1$.
If $B$ has at least two distinct prime divisors, there is no uniform bound
on when the increments of $S_{a,B}(n)$ become permanently equal to one.
\end{theorem}
\begin{proof}
Let $B=p^e$. The unit-component restriction forces $p\nmid a$.
For odd $p$, the order $d_p$ divides $p-1$, which divides $p^e-1=B-1$.
Hence the order condition is satisfied from the start, and
\[
S_{a,B}(n)=\left\lfloor\frac{ne+c_p}{e}\right\rfloor
=n+\left\lfloor\frac{c_p}{e}\right\rfloor.
\]
For $p=2$ the binary formula gives the same conclusion.

Now suppose $B$ has at least two prime divisors, and distinguish one of
them, $p_0$, with exponent $e_0$ in $B$. Choose any sufficiently large
integer $C\ge2$. By the Chinese remainder theorem choose $a\ge2$ satisfying
\begin{align*}
a&\equiv p_0\pmod{p_0^2},\\
a&\equiv1+q^{e_q C}\pmod{q^{e_q C+1}}
\quad\text{for each other prime }q\mid B.
\end{align*}
These moduli are pairwise coprime. The first condition gives
$\nu_{p_0}(a)=1$. Every other component is a unit congruent to $1$ modulo
$q$, with $c_q=e_qC$. At $q=2$, the requirement $e_qC\ge2$ ensures that
$\nu_2(a+1)=1$, so the same equality for $c_2$ holds.
There is no multiplicative-order obstruction. Formula~\eqref{eq:radix-envelope}
is therefore exact from $n=1$ and becomes
\begin{equation}\label{eq:crt-family}
S_{a,B}(n)=\min\left\{n+C,\left\lfloor\frac{B^n}{e_0}\right\rfloor\right\}.
\end{equation}
Fix any prescribed stage $n\ge2$ and take
$C\ge\lfloor B^n/e_0\rfloor$. At both $n-1$ and $n$ the exponential
term is the minimum. Their difference is
\[
\left\lfloor\frac{B^n}{e_0}\right\rfloor-
\left\lfloor\frac{B^{n-1}}{e_0}\right\rfloor>1.
\]
For completeness, $B\ge6$ and $e_0\le B$, so the difference of the
unrounded quantities is at least $B-1\ge5$ for $n\ge2$.
Yet eventually the linear term wins. Since the prescribed stage can be
arbitrarily late, there is no uniform onset bound.
\end{proof}

This isolates the cause of the decimal counterexample: simultaneous
components associated with different primes, some dividing the base and
others not. Extremely fast growth by itself is not the obstruction.

\section{What changes for genuine tetration?}\label{sec:tetration}

Define the integer tower sequence by
\[
T_0(a)=1,\qquad T_{h+1}(a)=a^{T_h(a)}.
\]
For $h\ge1$ this is $a\uparrow\uparrow h$. In contrast,
$X_b(a)=a^{10^b}$ satisfies $X_{b+1}(a)=X_b(a)^{10}$.
The two recurrences are fundamentally different.

Let $\Delta_h=T_{h+1}(a)-T_h(a)$. The tower recurrence yields
\begin{equation}\label{eq:tower-diff}
\Delta_h=a^{T_{h-1}(a)}\left(a^{\Delta_{h-1}}-1\right)
\qquad(h\ge2),
\end{equation}
and for $h=1$ the exponent in the second factor is $a-1$.
Local valuations are propagated through \emph{previous tower differences},
not through the predetermined exponent $9\cdot10^b$.

\subsection{A short derivation for base 3}
The following recovers a known frozen-digit formula in the tetration
literature~\cite{graham}; it is included to show the mechanism directly.

\begin{proposition}\label{prop:three}
For $T_h=T_h(3)$ and $h\ge1$,
\[
\nu_2(T_{h+1}-T_h)=2h+1,\qquad
\nu_5(T_{h+1}-T_h)=h-1.
\]
Consequently
\[
\nu_{10}(T_{h+1}-T_h)=h-1.
\]
\end{proposition}
\begin{proof}
The initial difference is $27-3=24$, with valuations $3$ and $0$.
Every difference is divisible by $4$. Since
$C_2(3)=2$ and $\nu_5(3^4-1)=1$, lifting applied to
\eqref{eq:tower-diff} increases the binary valuation by $2$ and the quinary
valuation by $1$ at each step. Induction gives both formulas.
\end{proof}
Because the decimal valuations strictly increase, the first difference in
a sum of later tower differences cannot be canceled at its first nonzero
digit. Thus $T_h(3)$ shares exactly $h-1$ trailing digits with every later
tower. This is an example where enormous growth and a simple exact digit
count coexist.

\subsection{The counterexample base behaves differently in a true tower}
\begin{proposition}\label{prop:odd-five-tower}
Suppose $a$ is an odd multiple of $5$. Put
$s=\nu_5(a)$, $c=C_2(a)$, and $A=\nu_2(a-1)$. For $h\ge1$,
\begin{equation}\label{eq:odd-five-tower}
\nu_{10}\!\left(T_{h+1}(a)-T_h(a)\right)
=\min\{sT_{h-1}(a),\ hc+A\}.
\end{equation}
\end{proposition}
\begin{proof}
For $h=1$, factor $a^a-a=a(a^{a-1}-1)$.
The quinary valuation is $s$, while binary lifting gives $c+A$.
At every subsequent stage, the second factor in~\eqref{eq:tower-diff}
is $-1$ modulo $5$, so the quinary valuation is $sT_{h-1}(a)$.
Every tower difference is even; binary lifting adds $c$ to its binary
valuation at each step. This gives $hc+A$ and proves the result.
\end{proof}
For $a=\astar$, we have $s=1$ and $A=c=99$, so the genuine tower profile is
\[
\min\{T_{h-1}(\astar),99(h+1)\}.
\]
At $h=1$ it is $1$, at $h=2$ it is $297$, and for every $h\ge2$ it equals
$99(h+1)$. To verify the last assertion, note $T_1(\astar)>297$ and
$T_h(\astar)$ increases at least geometrically, hence stays above the
indicated linear bound. Its subsequent increments are $99$, not $1$.
This does not conflict with the result about $\astar^{10^b}$; it emphasizes
why the two problems must be kept separate.

\section{Exact verification and reproducibility}\label{sec:verification}

\subsection{Certificates that do not expand the huge powers}
If a formula predicts that a difference has radix-$B$ valuation $s$, compute
it modulo $B^{s+1}$. A nonzero residue equal to $dB^s$ with
$1\le d<B$ proves that its valuation is exactly $s$.
Modular exponentiation performs this check without constructing the full
integer powers.

For the primary example, the independent certificates are
\begin{align}
\astar^{10^3}-\astar^{10^2}&\equiv4\cdot10^{100}\pmod{10^{101}},\label{eq:cert2}\\
\astar^{10^4}-\astar^{10^3}&\equiv5\cdot10^{102}\pmod{10^{103}}.\label{eq:cert3}
\end{align}
They give $S_{\astar}(2)=100$ and $S_{\astar}(3)=102$ directly.
The proofs earlier in the note do not depend on these calculations.

\subsection{Included programs and actual test results}
The archive contains \texttt{code/stabilization.py}, implementing the
formulas, and \texttt{code/verify.py}, which generates the result files.
Both use only the Python standard library and are written for Python~3.9
or later. The run packaged with this note used Python~3.13.5.

\begin{center}
\begin{tabular}{lr}
\toprule
Check category & Completed checks\\
\midrule
Primary-example modular certificates & 7\\
Small decimal bases, $2\le a\le2000$ & 8995\\
Delayed-family modular certificates & 64\\
Residue-class tests of the minimum table & 100\\
Stage-minimum modular certificates & 6\\
General-radix modular certificates, $2\le B\le35$ & 12489\\
Nonconsecutive decimal certificates & 432\\
Period-two radix-$14$ certificates & 4\\
Idempotent-limit tests & 20\\
\bottomrule
\end{tabular}
\end{center}
All completed checks passed. The general-radix sweep deliberately skipped
1111 cases whose predicted valuation exceeded a cap of 160, to avoid
constructing unnecessarily large verification moduli. No conclusion about
those skipped cases is drawn from the sweep; they are covered by the
mathematical theorem. The counts and exact certificates are preserved in
\texttt{data/verification.json} and \texttt{data/primary\_certificates.csv}.

The table test runs through $99\le K\le198$, covering every residue class
modulo $20$ repeatedly, and independently enumerates $j2^K\pm1$ for
$1\le j\le4$. Minimality over \emph{all} integers is established by
Theorem~\ref{thm:stage-minima}, not by that finite test.

\subsection{Reproduction commands and computational scope}
From the archive's root directory, run
\begin{lstlisting}
python code/verify.py
pdflatex -interaction=nonstopmode -halt-on-error delayed_digit_stabilization.tex
pdflatex -interaction=nonstopmode -halt-on-error delayed_digit_stabilization.tex
\end{lstlisting}
A PowerShell build script and a POSIX shell build script are also included.
The PDF uses standard TeX Live packages; no external bibliography tool is
required. The source contains its bibliography.

The modular verifier represents $B^n$ as an ordinary integer, so it is not
an evaluator for arbitrary unexpanded up-arrow expressions. The symbolic
minimum specification avoids expanding $2^K$ when a stage is too large.
Radix factorization in the utility code uses trial division and is intended
for small radices, not for cryptographic-size arbitrary integers. These
implementation limits do not restrict the proved formulas.

No proof-assistant formalization is included. The verification layer is
exact integer arithmetic plus independently computed residues, rather than
floating-point evidence or a formal kernel check.

\section{Conclusions and the remaining boundary}\label{sec:conclusions}

The selected conjecture has a negative answer. Its smallest counterexample
is a 31-digit integer with a compact description, $3\cdot2^{99}+1$.
The full explanation is the exact profile
\[
S_a(b)=\min\{t10^b,b+C\}
\]
for mixed decimal bases. The affine branch always wins eventually, but
there is no uniform bound on when it does so.

The result can be sharpened in several independent directions. The onset
is given by a single monotone integer inequality. The least exceptional
base at stage $b$ is explicit, with a coefficient and sign periodic in
$b$ modulo $20$. The exceptional set has an exact positive density, which
falls doubly exponentially with the stage. Arbitrary-radix analogues reveal
both multiplicative-order obstructions and a distinction between prime
power and mixed-prime radices. Fast-growing changes of index transform the
rate exactly as in~\eqref{eq:time-speed}.

These statements settle a modest, precisely posed open question and provide
a broader structure around it. They do not settle analytic interpolation
questions for tetration, transcendence questions about towers, or comparison
problems for arbitrary ordinal-indexed fast-growing hierarchies. Those are
different targets. Within the selected target, however, the answer is
complete: a counterexample, its minimality, an exact family of stage minima,
and quantitative descriptions of both delay and rarity.

\clearpage
\appendix
\section{A compact standalone answer to the original question}\label{app:short}
The following proof can be read without any other section.

\begin{resultbox}
The proposed threshold is false. Let $a=3\cdot2^{99}+1$.
Then $a\equiv15\pmod{25}$, so $\nu_5(a)=1$, while
$\nu_2(a-1)=99$ and $\nu_2(a+1)=1$.
For $b\ge2$,
\[
a^{10^{b+1}}-a^{10^b}=a^{10^b}\bigl(a^{9\cdot10^b}-1\bigr).
\]
Since $5\mid a$, its quinary valuation is $10^b$.
Binary lifting gives its binary valuation as
\[
99+1+b-1=b+99.
\]
Therefore $S_a(b)=\min\{10^b,b+99\}$, and
\[
D_a(3)=\min\{1000,102\}-\min\{100,101\}=2.
\]
Furthermore, for $a=5^K+1$ one obtains
$S_a(b)=\min\{10^b,b+K\}$. Choosing
$K=10^m-m+1$ gives $D_a(m+1)=2$ for every $m\ge2$.
Thus no fixed replacement threshold works for all bases.
\end{resultbox}

\section{A minimal executable certificate}\label{app:certificate}
This independent fragment uses no formulas from the companion module:
\begin{lstlisting}[language=Python]
a = 3 * 2**99 + 1

checks = [(2, 100, 4), (3, 102, 5)]
for b, s, expected_digit in checks:
    modulus = 10**(s + 1)
    residue = (pow(a, 10**(b + 1), modulus)
               - pow(a, 10**b, modulus)) % modulus
    assert residue == expected_digit * 10**s

assert 102 - 100 == 2
print("Counterexample verified exactly.")
\end{lstlisting}
The module's formula routine and the independent modular routine test
different descriptions of the same quantities. Agreement between them is
useful for catching transcription and indexing errors, while the proofs
establish the universal claims.

\clearpage
\phantomsection
\addcontentsline{toc}{section}{References}
\begin{thebibliography}{9}
\bibitem{question}
Marco Rip\`a.
\emph{Stabilization of one new trailing digit of $a^{10^b}$ as $b$ increases,
for every $a$ not divisible by $10$}.
Mathematics Stack Exchange, question 5103929, posted 24 October 2025.
\url{https://math.stackexchange.com/q/5103929}.
Public page checked for this note on 19 September 2026; the retrieved page
displayed no answers.

\bibitem{speed}
Marco Rip\`a.
The congruence speed formula.
\emph{Notes on Number Theory and Discrete Mathematics}
\textbf{27}(4) (2021), 43--61.
DOI: \url{https://doi.org/10.7546/nntdm.2021.27.4.43-61}.
Preprint: \url{https://arxiv.org/abs/2208.02622}.

\bibitem{stable}
Marco Rip\`a and Luca Onnis.
Number of stable digits of any integer tetration.
\emph{Notes on Number Theory and Discrete Mathematics}
\textbf{28}(3) (2022), 441--457.
DOI: \url{https://doi.org/10.7546/nntdm.2022.28.3.441-457}.
Preprint: \url{https://arxiv.org/abs/2210.07956}.

\bibitem{graham}
Marco Rip\`a.
Graham's number stable digits: An exact solution.
\emph{Notes on Number Theory and Discrete Mathematics}
\textbf{31}(3) (2025), 607--616.
Preprint version 2, 12 September 2025:
\url{https://arxiv.org/abs/2411.00015}.
\end{thebibliography}
\end{document}
